\documentclass[12pt, ukrainian]{article}
\usepackage[a4paper]{geometry}
\setlength\parindent{0mm}
%\setlength\parskip{4mm}
\geometry{top=1.3cm,bottom=1.3cm,left=1.5cm,right=1.5cm}
\pagestyle{empty}
\usepackage[T2A]{fontenc}
\usepackage[utf8]{inputenc}
\usepackage[ukrainian]{babel}
\begin{document}
%begin
{{\begin {scriptsize}\par
{ \center  Київський національний університет імені Тараса Шевченка \par}
{\parbox {5 cm}{Освітня програма \hfill}{\parbox {7 cm}{Системний аналіз\hfill}}\parbox {6 cm}{\hfill Семестр 2}}\par
{\parbox {5 cm} {Навчальний предмет \hfill}\parbox {7 cm}{Програмування \hfill}\parbox {6cm}{\hfill}\par}\vspace{-15pt} \end{scriptsize}}{\center Екзаменаційний білет \No \textbf { 1 }\par}}
{%beginitem
1. Математичний підхід до визначення складності обчислень.%enditem
\par
\vspace{2mm}

%beginitem
2. Кожен рядок текстового файлу містить по два дійсних числа, розділені білими символами.
Написати функцію, що для заданого текстового файлу будує новий текстовий файл, рядки якого крім зображень дйсних чисел ще містять результат обчислення на них функції $f$. У вихідному файлі зображення чисел мають розділятися пробілами. 
$$f(x, y) =\frac{\log_{10} x - 35}{(\arctg x - 0.6) (y - 11)}$$ 
Якість коду також оцінюється.

%enditem
\par
\vspace{2mm}

%beginitem
3. Реалізувати операцію злиття множин з повтореннями. Множина з повтореннями задається переліком своїх елементів, записаних за неспаданням у текстовому файлі через білі символи. Результат у такому ж самому форматі (і теж за неспаданням) має бути записаний у текстовий файл. Питома функція має отримувати три шляхи до файлів. Розмір вхідних множин такий, що вони не можуть бути завантажені повністю в пам’ять програми.

%enditem
}
{\smallskip\smallskip \begin {scriptsize} Затверджено на засіданні кафедри теоретичної кібернетики, протокол \No 4  від   25.01.2024 р.\par\smallskip\smallskip\smallskip\smallskip
{\parbox {6.5 cm} {Зав. кафедри \dotfill Ю.В.~Крак}}\hspace {0.7cm} {\hbox to 6.5 cm { Екзаменатор \dotfill Т.О.~Карнаух}} \end{scriptsize}}
%end
\vspace{0.9cm}
%begin
{{\begin {scriptsize}\par
{ \center  Київський національний університет імені Тараса Шевченка \par}
{\parbox {5 cm}{Освітня програма \hfill}{\parbox {7 cm}{Системний аналіз\hfill}}\parbox {6 cm}{\hfill Семестр 2}}\par
{\parbox {5 cm} {Навчальний предмет \hfill}\parbox {7 cm}{Програмування \hfill}\parbox {6cm}{\hfill}\par}\vspace{-15pt} \end{scriptsize}}{\center Екзаменаційний білет \No \textbf { 2 }\par}}
{%beginitem
1. Словники.%enditem
\par
\vspace{2mm}

%beginitem
2. Кожен рядок текстового файлу містить по два дійсних числа, розділені білими символами.
Написати функцію, що для заданого текстового файлу будує новий текстовий файл, рядки якого крім зображень дйсних чисел ще містять результат обчислення на них функції $f$. У вихідному файлі зображення чисел мають розділятися пробілами. 
$$f(x, y) =\frac{\sqrt x + 58}{(\cos x + 0.3) (y + 95)}$$ 
Якість коду також оцінюється.

%enditem
\par
\vspace{2mm}

%beginitem
3. Реалізувати операцію симетричної різниці множин з повтореннями. Множина з повтореннями задається переліком своїх елементів, записаних за неспаданням у текстовому файлі через білі символи. Результат у такому ж самому форматі (і теж за неспаданням) має бути записаний у текстовий файл. Питома функція має отримувати три шляхи до файлів. Розмір вхідних множин такий, що вони не можуть бути завантажені повністю в пам’ять програми.

%enditem
}
{\smallskip\smallskip \begin {scriptsize} Затверджено на засіданні кафедри теоретичної кібернетики, протокол \No 4  від   25.01.2024 р.\par\smallskip\smallskip\smallskip\smallskip
{\parbox {6.5 cm} {Зав. кафедри \dotfill Ю.В.~Крак}}\hspace {0.7cm} {\hbox to 6.5 cm { Екзаменатор \dotfill Т.О.~Карнаух}} \end{scriptsize}}
%end
\newpage
%begin
{{\begin {scriptsize}\par
{ \center  Київський національний університет імені Тараса Шевченка \par}
{\parbox {5 cm}{Освітня програма \hfill}{\parbox {7 cm}{Системний аналіз\hfill}}\parbox {6 cm}{\hfill Семестр 2}}\par
{\parbox {5 cm} {Навчальний предмет \hfill}\parbox {7 cm}{Програмування \hfill}\parbox {6cm}{\hfill}\par}\vspace{-15pt} \end{scriptsize}}{\center Екзаменаційний білет \No \textbf { 3 }\par}}
{%beginitem
1. Поняття часової та просторової складності обчислень. Складність у середньому. Асимптотичні оцінки складності.%enditem
\par
\vspace{2mm}

%beginitem
2. Кожен рядок текстового файлу містить по два дійсних числа, розділені білими символами.
Написати функцію, що для заданого текстового файлу будує новий текстовий файл, рядки якого крім зображень дйсних чисел ще містять результат обчислення на них функції $f$. У вихідному файлі зображення чисел мають розділятися пробілами. 
$$f(x, y) =\frac{\ln x - 15}{(\sin x - 0.8) (y - 61)}$$ 
Якість коду також оцінюється.

%enditem
\par
\vspace{2mm}

%beginitem
3. Реалізувати операцію перетину множин з повтореннями. Множина з повтореннями задається переліком своїх елементів, записаних за неспаданням у текстовому файлі через білі символи. Результат у такому ж самому форматі (і теж за неспаданням) має бути записаний у текстовий файл. Питома функція має отримувати три шляхи до файлів. Розмір вхідних множин такий, що вони не можуть бути завантажені повністю в пам’ять програми.

%enditem
}
{\smallskip\smallskip \begin {scriptsize} Затверджено на засіданні кафедри теоретичної кібернетики, протокол \No 4  від   25.01.2024 р.\par\smallskip\smallskip\smallskip\smallskip
{\parbox {6.5 cm} {Зав. кафедри \dotfill Ю.В.~Крак}}\hspace {0.7cm} {\hbox to 6.5 cm { Екзаменатор \dotfill Т.О.~Карнаух}} \end{scriptsize}}
%end
\vspace{0.9cm}
%begin
{{\begin {scriptsize}\par
{ \center  Київський національний університет імені Тараса Шевченка \par}
{\parbox {5 cm}{Освітня програма \hfill}{\parbox {7 cm}{Системний аналіз\hfill}}\parbox {6 cm}{\hfill Семестр 2}}\par
{\parbox {5 cm} {Навчальний предмет \hfill}\parbox {7 cm}{Програмування \hfill}\parbox {6cm}{\hfill}\par}\vspace{-15pt} \end{scriptsize}}{\center Екзаменаційний білет \No \textbf { 4 }\par}}
{%beginitem
1. Множини.%enditem
\par
\vspace{2mm}

%beginitem
2. Кожен рядок текстового файлу містить по два дійсних числа, розділені білими символами.
Написати функцію, що для заданого текстового файлу будує новий текстовий файл, рядки якого крім зображень дйсних чисел ще містять результат обчислення на них функції $f$. У вихідному файлі зображення чисел мають розділятися пробілами. 
$$f(x, y) =\frac{\log_{2} x + 13}{(\arccos x + 0.5) (y + 47)}$$ 
Якість коду також оцінюється.

%enditem
\par
\vspace{2mm}

%beginitem
3. Текстовий файл містить послідовність цілих чисел $a_1$, $a_2$, $\ldots$, $a_t$, записану у форматі одне число на рядок. Знайти сімнадцяте найбільше значення, що записане у файлі. Кожне значення, що записане у файлі, враховується стільки разів, скільки входить у файл. Кількість чисел у файлі така, що вони не можуть бути завантажені повністю в пам'ять програми. Розв'язання оформити у вигляді функції.

%enditem
}
{\smallskip\smallskip \begin {scriptsize} Затверджено на засіданні кафедри теоретичної кібернетики, протокол \No 4  від   25.01.2024 р.\par\smallskip\smallskip\smallskip\smallskip
{\parbox {6.5 cm} {Зав. кафедри \dotfill Ю.В.~Крак}}\hspace {0.7cm} {\hbox to 6.5 cm { Екзаменатор \dotfill Т.О.~Карнаух}} \end{scriptsize}}
%end
\newpage
%begin
{{\begin {scriptsize}\par
{ \center  Київський національний університет імені Тараса Шевченка \par}
{\parbox {5 cm}{Освітня програма \hfill}{\parbox {7 cm}{Системний аналіз\hfill}}\parbox {6 cm}{\hfill Семестр 2}}\par
{\parbox {5 cm} {Навчальний предмет \hfill}\parbox {7 cm}{Програмування \hfill}\parbox {6cm}{\hfill}\par}\vspace{-15pt} \end{scriptsize}}{\center Екзаменаційний білет \No \textbf { 5 }\par}}
{%beginitem
1. Емпіричний підхід до визначення складності обчислень.%enditem
\par
\vspace{2mm}

%beginitem
2. Кожен рядок текстового файлу містить по два дійсних числа, розділені білими символами.
Написати функцію, що для заданого текстового файлу будує новий текстовий файл, рядки якого крім зображень дйсних чисел ще містять результат обчислення на них функції $f$. У вихідному файлі зображення чисел мають розділятися пробілами. 
$$f(x, y) =\frac{\ x - 17}{(\tg x - 0.9) (y - 19)}$$ 
Якість коду також оцінюється.

%enditem
\par
\vspace{2mm}

%beginitem
3. Реалізувати операцію різниці множин з повтореннями. Множина з повтореннями задається переліком своїх елементів, записаних за неспаданням у текстовому файлі через білі символи. Результат у такому ж самому форматі (і теж за неспаданням) має бути записаний у текстовий файл. Питома функція має отримувати три шляхи до файлів. Розмір вхідних множин такий, що вони не можуть бути завантажені повністю в пам’ять програми.

%enditem
}
{\smallskip\smallskip \begin {scriptsize} Затверджено на засіданні кафедри теоретичної кібернетики, протокол \No 4  від   25.01.2024 р.\par\smallskip\smallskip\smallskip\smallskip
{\parbox {6.5 cm} {Зав. кафедри \dotfill Ю.В.~Крак}}\hspace {0.7cm} {\hbox to 6.5 cm { Екзаменатор \dotfill Т.О.~Карнаух}} \end{scriptsize}}
%end
\vspace{0.9cm}
%begin
{{\begin {scriptsize}\par
{ \center  Київський національний університет імені Тараса Шевченка \par}
{\parbox {5 cm}{Освітня програма \hfill}{\parbox {7 cm}{Системний аналіз\hfill}}\parbox {6 cm}{\hfill Семестр 2}}\par
{\parbox {5 cm} {Навчальний предмет \hfill}\parbox {7 cm}{Програмування \hfill}\parbox {6cm}{\hfill}\par}\vspace{-15pt} \end{scriptsize}}{\center Екзаменаційний білет \No \textbf { 6 }\par}}
{%beginitem
1. Обробка файлів: режими відкриття файлів.%enditem
\par
\vspace{2mm}

%beginitem
2. Кожен рядок текстового файлу містить по два дійсних числа, розділені білими символами.
Написати функцію, що для заданого текстового файлу будує новий текстовий файл, рядки якого крім зображень дйсних чисел ще містять результат обчислення на них функції $f$. У вихідному файлі зображення чисел мають розділятися пробілами. 
$$f(x, y) =\frac{\pi x + 36}{(\arcsin x + 0.4) (y + 69)}$$ 
Якість коду також оцінюється.

%enditem
\par
\vspace{2mm}

%beginitem
3. Текстовий файл містить послідовність цілих чисел $a_1$, $a_2$, $\ldots$, $a_t$, записану у форматі одне число на рядок. Знайти тринадцяте найменше значення, що записане у файлі. Кожне значення, що записане у файлі, враховується тільки один раз. Кількість чисел у файлі така, що вони не можуть бути завантажені повністю в пам'ять програми. Розв'язання оформити у вигляді функції.

%enditem
}
{\smallskip\smallskip \begin {scriptsize} Затверджено на засіданні кафедри теоретичної кібернетики, протокол \No 4  від   25.01.2024 р.\par\smallskip\smallskip\smallskip\smallskip
{\parbox {6.5 cm} {Зав. кафедри \dotfill Ю.В.~Крак}}\hspace {0.7cm} {\hbox to 6.5 cm { Екзаменатор \dotfill Т.О.~Карнаух}} \end{scriptsize}}
%end
\newpage
%begin
{{\begin {scriptsize}\par
{ \center  Київський національний університет імені Тараса Шевченка \par}
{\parbox {5 cm}{Освітня програма \hfill}{\parbox {7 cm}{Системний аналіз\hfill}}\parbox {6 cm}{\hfill Семестр 2}}\par
{\parbox {5 cm} {Навчальний предмет \hfill}\parbox {7 cm}{Програмування \hfill}\parbox {6cm}{\hfill}\par}\vspace{-15pt} \end{scriptsize}}{\center Екзаменаційний білет \No \textbf { 7 }\par}}
{%beginitem
1. Файли формату csv.%enditem
\par
\vspace{2mm}

%beginitem
2. Кожен рядок текстового файлу містить по два дійсних числа, розділені білими символами.
Написати функцію, що для заданого текстового файлу будує новий текстовий файл, рядки якого крім зображень дйсних чисел ще містять результат обчислення на них функції $f$. У вихідному файлі зображення чисел мають розділятися пробілами. 
$$f(x, y) =\frac{\ x - 63}{(\ctg x + 0.7) (y - 53)}$$ 
Якість коду також оцінюється.

%enditem
\par
\vspace{2mm}

%beginitem
3. Текстовий файл містить послідовність цілих чисел $a_1$, $a_2$, $\ldots$, $a_t$, записану у форматі одне число на рядок. Знайти сімнадцяте найменше значення, що записане у файлі. Кожне значення, що записане у файлі, враховується стільки разів, скільки входить у файл. Кількість чисел у файлі така, що вони не можуть бути завантажені повністю в пам'ять програми. Розв'язання оформити у вигляді функції.

%enditem
}
{\smallskip\smallskip \begin {scriptsize} Затверджено на засіданні кафедри теоретичної кібернетики, протокол \No 4  від   25.01.2024 р.\par\smallskip\smallskip\smallskip\smallskip
{\parbox {6.5 cm} {Зав. кафедри \dotfill Ю.В.~Крак}}\hspace {0.7cm} {\hbox to 6.5 cm { Екзаменатор \dotfill Т.О.~Карнаух}} \end{scriptsize}}
%end
\vspace{0.9cm}
%begin
{{\begin {scriptsize}\par
{ \center  Київський національний університет імені Тараса Шевченка \par}
{\parbox {5 cm}{Освітня програма \hfill}{\parbox {7 cm}{Системний аналіз\hfill}}\parbox {6 cm}{\hfill Семестр 2}}\par
{\parbox {5 cm} {Навчальний предмет \hfill}\parbox {7 cm}{Програмування \hfill}\parbox {6cm}{\hfill}\par}\vspace{-15pt} \end{scriptsize}}{\center Екзаменаційний білет \No \textbf { 8 }\par}}
{%beginitem
1. Алгоритм пошуку в ширину та задачі, які він може вирішувати.%enditem
\par
\vspace{2mm}

%beginitem
2. Кожен рядок текстового файлу містить по два дійсних числа, розділені білими символами.
Написати функцію, що для заданого текстового файлу будує новий текстовий файл, рядки якого крім зображень дйсних чисел ще містять результат обчислення на них функції $f$. У вихідному файлі зображення чисел мають розділятися пробілами. 
$$f(x, y) =\frac{\pi x + 28}{(\cos x - 0.1) (y + 60)}$$ 
Якість коду також оцінюється.

%enditem
\par
\vspace{2mm}

%beginitem
3. Текстовий файл містить послідовність цілих чисел $a_1$, $a_2$, $\ldots$, $a_t$, записану у форматі одне число на рядок. Знайти тринадцяте найбільше значення, що записане у файлі. Кожне значення, що записане у файлі, враховується тільки один раз. Кількість чисел у файлі така, що вони не можуть бути завантажені повністю в пам'ять програми. Розв'язання оформити у вигляді функції.

%enditem
}
{\smallskip\smallskip \begin {scriptsize} Затверджено на засіданні кафедри теоретичної кібернетики, протокол \No 4  від   25.01.2024 р.\par\smallskip\smallskip\smallskip\smallskip
{\parbox {6.5 cm} {Зав. кафедри \dotfill Ю.В.~Крак}}\hspace {0.7cm} {\hbox to 6.5 cm { Екзаменатор \dotfill Т.О.~Карнаух}} \end{scriptsize}}
%end
\newpage
%begin
{{\begin {scriptsize}\par
{ \center  Київський національний університет імені Тараса Шевченка \par}
{\parbox {5 cm}{Освітня програма \hfill}{\parbox {7 cm}{Системний аналіз\hfill}}\parbox {6 cm}{\hfill Семестр 2}}\par
{\parbox {5 cm} {Навчальний предмет \hfill}\parbox {7 cm}{Програмування \hfill}\parbox {6cm}{\hfill}\par}\vspace{-15pt} \end{scriptsize}}{\center Екзаменаційний білет \No \textbf { 9 }\par}}
{%beginitem
1. Багатовимірні масиви бібліотеки numpy: у чому відмінність моделювання матриці багатовимірним масивом бібліотеки numpy та як список списків.%enditem
\par
\vspace{2mm}

%beginitem
2. Кожен рядок текстового файлу містить по два дійсних числа, розділені білими символами.
Написати функцію, що для заданого текстового файлу будує новий текстовий файл, рядки якого крім зображень дйсних чисел ще містять результат обчислення на них функції $f$. У вихідному файлі зображення чисел мають розділятися пробілами. 
$$f(x, y) =\frac{\sqrt x + 55}{(\arccos x + 0.2) (y - 64)}$$ 
Якість коду також оцінюється.

%enditem
\par
\vspace{2mm}

%beginitem
3. Кожен з двох текстових файлів містить послідовність дійсних чисел, записану у форматі одне число на рядок.  Перевірити, чи задають вони одну ту саму послідовність. Кількість чисел у кожному з двох файлів така, що вони не можуть бути завантажені повністю в пам'ять програми. Розв'язання оформити у вигляді функції.

%enditem
}
{\smallskip\smallskip \begin {scriptsize} Затверджено на засіданні кафедри теоретичної кібернетики, протокол \No 4  від   25.01.2024 р.\par\smallskip\smallskip\smallskip\smallskip
{\parbox {6.5 cm} {Зав. кафедри \dotfill Ю.В.~Крак}}\hspace {0.7cm} {\hbox to 6.5 cm { Екзаменатор \dotfill Т.О.~Карнаух}} \end{scriptsize}}
%end
\vspace{0.9cm}
%begin
{{\begin {scriptsize}\par
{ \center  Київський національний університет імені Тараса Шевченка \par}
{\parbox {5 cm}{Освітня програма \hfill}{\parbox {7 cm}{Системний аналіз\hfill}}\parbox {6 cm}{\hfill Семестр 2}}\par
{\parbox {5 cm} {Навчальний предмет \hfill}\parbox {7 cm}{Програмування \hfill}\parbox {6cm}{\hfill}\par}\vspace{-15pt} \end{scriptsize}}{\center Екзаменаційний білет \No \textbf { 10 }\par}}
{%beginitem
1. Лінійні зв’язані структури даних.%enditem
\par
\vspace{2mm}

%beginitem
2. Кожен рядок текстового файлу містить по два дійсних числа, розділені білими символами.
Написати функцію, що для заданого текстового файлу будує новий текстовий файл, рядки якого крім зображень дйсних чисел ще містять результат обчислення на них функції $f$. У вихідному файлі зображення чисел мають розділятися пробілами. 
$$f(x, y) =\frac{\log_{2} x - 72}{(\sin x - 0.9) (y + 37)}$$ 
Якість коду також оцінюється.

%enditem
\par
\vspace{2mm}

%beginitem
3. Текстовий файл містить послідовність цілих чисел $a_1$, $a_2$, $\ldots$, $a_t$, записану у форматі одне число на рядок. Знайти тринадцяте найбільше значення, що записане у файлі. Кожне значення, що записане у файлі, враховується тільки один раз. Кількість чисел у файлі така, що вони не можуть бути завантажені повністю в пам'ять програми. Розв'язання оформити у вигляді функції.

%enditem
}
{\smallskip\smallskip \begin {scriptsize} Затверджено на засіданні кафедри теоретичної кібернетики, протокол \No 4  від   25.01.2024 р.\par\smallskip\smallskip\smallskip\smallskip
{\parbox {6.5 cm} {Зав. кафедри \dotfill Ю.В.~Крак}}\hspace {0.7cm} {\hbox to 6.5 cm { Екзаменатор \dotfill Т.О.~Карнаух}} \end{scriptsize}}
%end
\newpage
%begin
{{\begin {scriptsize}\par
{ \center  Київський національний університет імені Тараса Шевченка \par}
{\parbox {5 cm}{Освітня програма \hfill}{\parbox {7 cm}{Системний аналіз\hfill}}\parbox {6 cm}{\hfill Семестр 2}}\par
{\parbox {5 cm} {Навчальний предмет \hfill}\parbox {7 cm}{Програмування \hfill}\parbox {6cm}{\hfill}\par}\vspace{-15pt} \end{scriptsize}}{\center Екзаменаційний білет \No \textbf { 11 }\par}}
{%beginitem
1. Поняття стеку та черги. Основні операції зі стеками та чергами.%enditem
\par
\vspace{2mm}

%beginitem
2. Кожен рядок текстового файлу містить по два дійсних числа, розділені білими символами.
Написати функцію, що для заданого текстового файлу будує новий текстовий файл, рядки якого крім зображень дйсних чисел ще містять результат обчислення на них функції $f$. У вихідному файлі зображення чисел мають розділятися пробілами. 
$$f(x, y) =\frac{\log_{10} x + 23}{(\arcsin x + 0.4) (y + 52)}$$ 
Якість коду також оцінюється.

%enditem
\par
\vspace{2mm}

%beginitem
3. Реалізувати операцію симетричної різниці множин з повтореннями. Множина з повтореннями задається переліком своїх елементів, записаних за неспаданням у текстовому файлі через білі символи. Результат у такому ж самому форматі (і теж за неспаданням) має бути записаний у текстовий файл. Питома функція має отримувати три шляхи до файлів. Розмір вхідних множин такий, що вони не можуть бути завантажені повністю в пам’ять програми.

%enditem
}
{\smallskip\smallskip \begin {scriptsize} Затверджено на засіданні кафедри теоретичної кібернетики, протокол \No 4  від   25.01.2024 р.\par\smallskip\smallskip\smallskip\smallskip
{\parbox {6.5 cm} {Зав. кафедри \dotfill Ю.В.~Крак}}\hspace {0.7cm} {\hbox to 6.5 cm { Екзаменатор \dotfill Т.О.~Карнаух}} \end{scriptsize}}
%end
\vspace{0.9cm}
%begin
{{\begin {scriptsize}\par
{ \center  Київський національний університет імені Тараса Шевченка \par}
{\parbox {5 cm}{Освітня програма \hfill}{\parbox {7 cm}{Системний аналіз\hfill}}\parbox {6 cm}{\hfill Семестр 2}}\par
{\parbox {5 cm} {Навчальний предмет \hfill}\parbox {7 cm}{Програмування \hfill}\parbox {6cm}{\hfill}\par}\vspace{-15pt} \end{scriptsize}}{\center Екзаменаційний білет \No \textbf { 12 }\par}}
{%beginitem
1. Дерева. Обходи дерев.%enditem
\par
\vspace{2mm}

%beginitem
2. Кожен рядок текстового файлу містить по два дійсних числа, розділені білими символами.
Написати функцію, що для заданого текстового файлу будує новий текстовий файл, рядки якого крім зображень дйсних чисел ще містять результат обчислення на них функції $f$. У вихідному файлі зображення чисел мають розділятися пробілами. 
$$f(x, y) =\frac{\ln x - 66}{(\ctg x - 0.5) (y - 46)}$$ 
Якість коду також оцінюється.

%enditem
\par
\vspace{2mm}

%beginitem
3. Реалізувати операцію злиття множин з повтореннями. Множина з повтореннями задається переліком своїх елементів, записаних за неспаданням у текстовому файлі через білі символи. Результат у такому ж самому форматі (і теж за неспаданням) має бути записаний у текстовий файл. Питома функція має отримувати три шляхи до файлів. Розмір вхідних множин такий, що вони не можуть бути завантажені повністю в пам’ять програми.

%enditem
}
{\smallskip\smallskip \begin {scriptsize} Затверджено на засіданні кафедри теоретичної кібернетики, протокол \No 4  від   25.01.2024 р.\par\smallskip\smallskip\smallskip\smallskip
{\parbox {6.5 cm} {Зав. кафедри \dotfill Ю.В.~Крак}}\hspace {0.7cm} {\hbox to 6.5 cm { Екзаменатор \dotfill Т.О.~Карнаух}} \end{scriptsize}}
%end
\newpage
%begin
{{\begin {scriptsize}\par
{ \center  Київський національний університет імені Тараса Шевченка \par}
{\parbox {5 cm}{Освітня програма \hfill}{\parbox {7 cm}{Системний аналіз\hfill}}\parbox {6 cm}{\hfill Семестр 2}}\par
{\parbox {5 cm} {Навчальний предмет \hfill}\parbox {7 cm}{Програмування \hfill}\parbox {6cm}{\hfill}\par}\vspace{-15pt} \end{scriptsize}}{\center Екзаменаційний білет \No \textbf { 13 }\par}}
{%beginitem
1. Ітератори.%enditem
\par
\vspace{2mm}

%beginitem
2. Кожен рядок текстового файлу містить по два дійсних числа, розділені білими символами.
Написати функцію, що для заданого текстового файлу будує новий текстовий файл, рядки якого крім зображень дйсних чисел ще містять результат обчислення на них функції $f$. У вихідному файлі зображення чисел мають розділятися пробілами. 
$$f(x, y) =\frac{\sqrt x + 62}{(\arctg x - 0.2) (y - 90)}$$ 
Якість коду також оцінюється.

%enditem
\par
\vspace{2mm}

%beginitem
3. Кожен з двох текстових файлів містить послідовність дійсних чисел, записану у форматі одне число на рядок.  Перевірити, чи задають вони одну ту саму послідовність. Кількість чисел у кожному з двох файлів така, що вони не можуть бути завантажені повністю в пам'ять програми. Розв'язання оформити у вигляді функції.

%enditem
}
{\smallskip\smallskip \begin {scriptsize} Затверджено на засіданні кафедри теоретичної кібернетики, протокол \No 4  від   25.01.2024 р.\par\smallskip\smallskip\smallskip\smallskip
{\parbox {6.5 cm} {Зав. кафедри \dotfill Ю.В.~Крак}}\hspace {0.7cm} {\hbox to 6.5 cm { Екзаменатор \dotfill Т.О.~Карнаух}} \end{scriptsize}}
%end
\vspace{0.9cm}
%begin
{{\begin {scriptsize}\par
{ \center  Київський національний університет імені Тараса Шевченка \par}
{\parbox {5 cm}{Освітня програма \hfill}{\parbox {7 cm}{Системний аналіз\hfill}}\parbox {6 cm}{\hfill Семестр 2}}\par
{\parbox {5 cm} {Навчальний предмет \hfill}\parbox {7 cm}{Програмування \hfill}\parbox {6cm}{\hfill}\par}\vspace{-15pt} \end{scriptsize}}{\center Екзаменаційний білет \No \textbf { 14 }\par}}
{%beginitem
1. Статичні поля, статичні методи, методи класів.%enditem
\par
\vspace{2mm}

%beginitem
2. Кожен рядок текстового файлу містить по два дійсних числа, розділені білими символами.
Написати функцію, що для заданого текстового файлу будує новий текстовий файл, рядки якого крім зображень дйсних чисел ще містять результат обчислення на них функції $f$. У вихідному файлі зображення чисел мають розділятися пробілами. 
$$f(x, y) =\frac{\ln x - 65}{(\tg x + 0.3) (y + 68)}$$ 
Якість коду також оцінюється.

%enditem
\par
\vspace{2mm}

%beginitem
3. Текстовий файл містить послідовність цілих чисел $a_1$, $a_2$, $\ldots$, $a_t$, записану у форматі одне число на рядок. Знайти тринадцяте найменше значення, що записане у файлі. Кожне значення, що записане у файлі, враховується тільки один раз. Кількість чисел у файлі така, що вони не можуть бути завантажені повністю в пам'ять програми. Розв'язання оформити у вигляді функції.

%enditem
}
{\smallskip\smallskip \begin {scriptsize} Затверджено на засіданні кафедри теоретичної кібернетики, протокол \No 4  від   25.01.2024 р.\par\smallskip\smallskip\smallskip\smallskip
{\parbox {6.5 cm} {Зав. кафедри \dotfill Ю.В.~Крак}}\hspace {0.7cm} {\hbox to 6.5 cm { Екзаменатор \dotfill Т.О.~Карнаух}} \end{scriptsize}}
%end
\newpage
%begin
{{\begin {scriptsize}\par
{ \center  Київський національний університет імені Тараса Шевченка \par}
{\parbox {5 cm}{Освітня програма \hfill}{\parbox {7 cm}{Системний аналіз\hfill}}\parbox {6 cm}{\hfill Семестр 2}}\par
{\parbox {5 cm} {Навчальний предмет \hfill}\parbox {7 cm}{Програмування \hfill}\parbox {6cm}{\hfill}\par}\vspace{-15pt} \end{scriptsize}}{\center Екзаменаційний білет \No \textbf { 15 }\par}}
{%beginitem
1. Контекстні менеджери.%enditem
\par
\vspace{2mm}

%beginitem
2. Кожен рядок текстового файлу містить по два дійсних числа, розділені білими символами.
Написати функцію, що для заданого текстового файлу будує новий текстовий файл, рядки якого крім зображень дйсних чисел ще містять результат обчислення на них функції $f$. У вихідному файлі зображення чисел мають розділятися пробілами. 
$$f(x, y) =\frac{\pi x - 51}{(\sin x - 0.1) (y - 30)}$$ 
Якість коду також оцінюється.

%enditem
\par
\vspace{2mm}

%beginitem
3. Текстовий файл містить послідовність цілих чисел $a_1$, $a_2$, $\ldots$, $a_t$, записану у форматі одне число на рядок. Знайти сімнадцяте найбільше значення, що записане у файлі. Кожне значення, що записане у файлі, враховується стільки разів, скільки входить у файл. Кількість чисел у файлі така, що вони не можуть бути завантажені повністю в пам'ять програми. Розв'язання оформити у вигляді функції.

%enditem
}
{\smallskip\smallskip \begin {scriptsize} Затверджено на засіданні кафедри теоретичної кібернетики, протокол \No 4  від   25.01.2024 р.\par\smallskip\smallskip\smallskip\smallskip
{\parbox {6.5 cm} {Зав. кафедри \dotfill Ю.В.~Крак}}\hspace {0.7cm} {\hbox to 6.5 cm { Екзаменатор \dotfill Т.О.~Карнаух}} \end{scriptsize}}
%end
\vspace{0.9cm}
%begin
{{\begin {scriptsize}\par
{ \center  Київський національний університет імені Тараса Шевченка \par}
{\parbox {5 cm}{Освітня програма \hfill}{\parbox {7 cm}{Системний аналіз\hfill}}\parbox {6 cm}{\hfill Семестр 2}}\par
{\parbox {5 cm} {Навчальний предмет \hfill}\parbox {7 cm}{Програмування \hfill}\parbox {6cm}{\hfill}\par}\vspace{-15pt} \end{scriptsize}}{\center Екзаменаційний білет \No \textbf { 16 }\par}}
{%beginitem
1. Файли формату json.%enditem
\par
\vspace{2mm}

%beginitem
2. Кожен рядок текстового файлу містить по два дійсних числа, розділені білими символами.
Написати функцію, що для заданого текстового файлу будує новий текстовий файл, рядки якого крім зображень дйсних чисел ще містять результат обчислення на них функції $f$. У вихідному файлі зображення чисел мають розділятися пробілами. 
$$f(x, y) =\frac{\log_{10} x + 96}{(\arccos x + 0.7) (y + 99)}$$ 
Якість коду також оцінюється.

%enditem
\par
\vspace{2mm}

%beginitem
3. Текстовий файл містить послідовність цілих чисел $a_1$, $a_2$, $\ldots$, $a_t$, записану у форматі одне число на рядок. Знайти сімнадцяте найменше значення, що записане у файлі. Кожне значення, що записане у файлі, враховується стільки разів, скільки входить у файл. Кількість чисел у файлі така, що вони не можуть бути завантажені повністю в пам'ять програми. Розв'язання оформити у вигляді функції.

%enditem
}
{\smallskip\smallskip \begin {scriptsize} Затверджено на засіданні кафедри теоретичної кібернетики, протокол \No 4  від   25.01.2024 р.\par\smallskip\smallskip\smallskip\smallskip
{\parbox {6.5 cm} {Зав. кафедри \dotfill Ю.В.~Крак}}\hspace {0.7cm} {\hbox to 6.5 cm { Екзаменатор \dotfill Т.О.~Карнаух}} \end{scriptsize}}
%end
\newpage
%begin
{{\begin {scriptsize}\par
{ \center  Київський національний університет імені Тараса Шевченка \par}
{\parbox {5 cm}{Освітня програма \hfill}{\parbox {7 cm}{Системний аналіз\hfill}}\parbox {6 cm}{\hfill Семестр 2}}\par
{\parbox {5 cm} {Навчальний предмет \hfill}\parbox {7 cm}{Програмування \hfill}\parbox {6cm}{\hfill}\par}\vspace{-15pt} \end{scriptsize}}{\center Екзаменаційний білет \No \textbf { 17 }\par}}
{%beginitem
1. Алгоритм пошуку в глибину та задачі, які він може вирішувати.%enditem
\par
\vspace{2mm}

%beginitem
2. Кожен рядок текстового файлу містить по два дійсних числа, розділені білими символами.
Написати функцію, що для заданого текстового файлу будує новий текстовий файл, рядки якого крім зображень дйсних чисел ще містять результат обчислення на них функції $f$. У вихідному файлі зображення чисел мають розділятися пробілами. 
$$f(x, y) =\frac{\ x + 25}{(\tg x - 0.8) (y - 84)}$$ 
Якість коду також оцінюється.

%enditem
\par
\vspace{2mm}

%beginitem
3. Реалізувати операцію перетину множин з повтореннями. Множина з повтореннями задається переліком своїх елементів, записаних за неспаданням у текстовому файлі через білі символи. Результат у такому ж самому форматі (і теж за неспаданням) має бути записаний у текстовий файл. Питома функція має отримувати три шляхи до файлів. Розмір вхідних множин такий, що вони не можуть бути завантажені повністю в пам’ять програми.

%enditem
}
{\smallskip\smallskip \begin {scriptsize} Затверджено на засіданні кафедри теоретичної кібернетики, протокол \No 4  від   25.01.2024 р.\par\smallskip\smallskip\smallskip\smallskip
{\parbox {6.5 cm} {Зав. кафедри \dotfill Ю.В.~Крак}}\hspace {0.7cm} {\hbox to 6.5 cm { Екзаменатор \dotfill Т.О.~Карнаух}} \end{scriptsize}}
%end
\vspace{0.9cm}
%begin
{{\begin {scriptsize}\par
{ \center  Київський національний університет імені Тараса Шевченка \par}
{\parbox {5 cm}{Освітня програма \hfill}{\parbox {7 cm}{Системний аналіз\hfill}}\parbox {6 cm}{\hfill Семестр 2}}\par
{\parbox {5 cm} {Навчальний предмет \hfill}\parbox {7 cm}{Програмування \hfill}\parbox {6cm}{\hfill}\par}\vspace{-15pt} \end{scriptsize}}{\center Екзаменаційний білет \No \textbf { 18 }\par}}
{%beginitem
1. Контекстні менеджери.%enditem
\par
\vspace{2mm}

%beginitem
2. Кожен рядок текстового файлу містить по два дійсних числа, розділені білими символами.
Написати функцію, що для заданого текстового файлу будує новий текстовий файл, рядки якого крім зображень дйсних чисел ще містять результат обчислення на них функції $f$. У вихідному файлі зображення чисел мають розділятися пробілами. 
$$f(x, y) =\frac{\log_{2} x - 59}{(\ctg x + 0.6) (y + 54)}$$ 
Якість коду також оцінюється.

%enditem
\par
\vspace{2mm}

%beginitem
3. Реалізувати операцію різниці множин з повтореннями. Множина з повтореннями задається переліком своїх елементів, записаних за неспаданням у текстовому файлі через білі символи. Результат у такому ж самому форматі (і теж за неспаданням) має бути записаний у текстовий файл. Питома функція має отримувати три шляхи до файлів. Розмір вхідних множин такий, що вони не можуть бути завантажені повністю в пам’ять програми.

%enditem
}
{\smallskip\smallskip \begin {scriptsize} Затверджено на засіданні кафедри теоретичної кібернетики, протокол \No 4  від   25.01.2024 р.\par\smallskip\smallskip\smallskip\smallskip
{\parbox {6.5 cm} {Зав. кафедри \dotfill Ю.В.~Крак}}\hspace {0.7cm} {\hbox to 6.5 cm { Екзаменатор \dotfill Т.О.~Карнаух}} \end{scriptsize}}
%end
\newpage
%begin
{{\begin {scriptsize}\par
{ \center  Київський національний університет імені Тараса Шевченка \par}
{\parbox {5 cm}{Освітня програма \hfill}{\parbox {7 cm}{Системний аналіз\hfill}}\parbox {6 cm}{\hfill Семестр 2}}\par
{\parbox {5 cm} {Навчальний предмет \hfill}\parbox {7 cm}{Програмування \hfill}\parbox {6cm}{\hfill}\par}\vspace{-15pt} \end{scriptsize}}{\center Екзаменаційний білет \No \textbf { 19 }\par}}
{%beginitem
1. Файли формату json.%enditem
\par
\vspace{2mm}

%beginitem
2. Кожен рядок текстового файлу містить по два дійсних числа, розділені білими символами.
Написати функцію, що для заданого текстового файлу будує новий текстовий файл, рядки якого крім зображень дйсних чисел ще містять результат обчислення на них функції $f$. У вихідному файлі зображення чисел мають розділятися пробілами. 
$$f(x, y) =\frac{\ x + 79}{(\arcsin x - 0.4) (y + 44)}$$ 
Якість коду також оцінюється.

%enditem
\par
\vspace{2mm}

%beginitem
3. Реалізувати операцію симетричної різниці множин з повтореннями. Множина з повтореннями задається переліком своїх елементів, записаних за неспаданням у текстовому файлі через білі символи. Результат у такому ж самому форматі (і теж за неспаданням) має бути записаний у текстовий файл. Питома функція має отримувати три шляхи до файлів. Розмір вхідних множин такий, що вони не можуть бути завантажені повністю в пам’ять програми.

%enditem
}
{\smallskip\smallskip \begin {scriptsize} Затверджено на засіданні кафедри теоретичної кібернетики, протокол \No 4  від   25.01.2024 р.\par\smallskip\smallskip\smallskip\smallskip
{\parbox {6.5 cm} {Зав. кафедри \dotfill Ю.В.~Крак}}\hspace {0.7cm} {\hbox to 6.5 cm { Екзаменатор \dotfill Т.О.~Карнаух}} \end{scriptsize}}
%end
\vspace{0.9cm}
%begin
{{\begin {scriptsize}\par
{ \center  Київський національний університет імені Тараса Шевченка \par}
{\parbox {5 cm}{Освітня програма \hfill}{\parbox {7 cm}{Системний аналіз\hfill}}\parbox {6 cm}{\hfill Семестр 2}}\par
{\parbox {5 cm} {Навчальний предмет \hfill}\parbox {7 cm}{Програмування \hfill}\parbox {6cm}{\hfill}\par}\vspace{-15pt} \end{scriptsize}}{\center Екзаменаційний білет \No \textbf { 20 }\par}}
{%beginitem
1. Поняття стеку та черги. Основні операції зі стеками та чергами.%enditem
\par
\vspace{2mm}

%beginitem
2. Кожен рядок текстового файлу містить по два дійсних числа, розділені білими символами.
Написати функцію, що для заданого текстового файлу будує новий текстовий файл, рядки якого крім зображень дйсних чисел ще містять результат обчислення на них функції $f$. У вихідному файлі зображення чисел мають розділятися пробілами. 
$$f(x, y) =\frac{\log_{2} x - 94}{(\cos x + 0.9) (y - 82)}$$ 
Якість коду також оцінюється.

%enditem
\par
\vspace{2mm}

%beginitem
3. Текстовий файл містить послідовність цілих чисел $a_1$, $a_2$, $\ldots$, $a_t$, записану у форматі одне число на рядок. Знайти сімнадцяте найбільше значення, що записане у файлі. Кожне значення, що записане у файлі, враховується стільки разів, скільки входить у файл. Кількість чисел у файлі така, що вони не можуть бути завантажені повністю в пам'ять програми. Розв'язання оформити у вигляді функції.

%enditem
}
{\smallskip\smallskip \begin {scriptsize} Затверджено на засіданні кафедри теоретичної кібернетики, протокол \No 4  від   25.01.2024 р.\par\smallskip\smallskip\smallskip\smallskip
{\parbox {6.5 cm} {Зав. кафедри \dotfill Ю.В.~Крак}}\hspace {0.7cm} {\hbox to 6.5 cm { Екзаменатор \dotfill Т.О.~Карнаух}} \end{scriptsize}}
%end
\newpage
%begin
{{\begin {scriptsize}\par
{ \center  Київський національний університет імені Тараса Шевченка \par}
{\parbox {5 cm}{Освітня програма \hfill}{\parbox {7 cm}{Системний аналіз\hfill}}\parbox {6 cm}{\hfill Семестр 2}}\par
{\parbox {5 cm} {Навчальний предмет \hfill}\parbox {7 cm}{Програмування \hfill}\parbox {6cm}{\hfill}\par}\vspace{-15pt} \end{scriptsize}}{\center Екзаменаційний білет \No \textbf { 21 }\par}}
{%beginitem
1. Обробка файлів: режими відкриття файлів.%enditem
\par
\vspace{2mm}

%beginitem
2. Кожен рядок текстового файлу містить по два дійсних числа, розділені білими символами.
Написати функцію, що для заданого текстового файлу будує новий текстовий файл, рядки якого крім зображень дйсних чисел ще містять результат обчислення на них функції $f$. У вихідному файлі зображення чисел мають розділятися пробілами. 
$$f(x, y) =\frac{\sqrt x - 33}{(\arctg x + 0.1) (y + 85)}$$ 
Якість коду також оцінюється.

%enditem
\par
\vspace{2mm}

%beginitem
3. Реалізувати операцію злиття множин з повтореннями. Множина з повтореннями задається переліком своїх елементів, записаних за неспаданням у текстовому файлі через білі символи. Результат у такому ж самому форматі (і теж за неспаданням) має бути записаний у текстовий файл. Питома функція має отримувати три шляхи до файлів. Розмір вхідних множин такий, що вони не можуть бути завантажені повністю в пам’ять програми.

%enditem
}
{\smallskip\smallskip \begin {scriptsize} Затверджено на засіданні кафедри теоретичної кібернетики, протокол \No 4  від   25.01.2024 р.\par\smallskip\smallskip\smallskip\smallskip
{\parbox {6.5 cm} {Зав. кафедри \dotfill Ю.В.~Крак}}\hspace {0.7cm} {\hbox to 6.5 cm { Екзаменатор \dotfill Т.О.~Карнаух}} \end{scriptsize}}
%end
\vspace{0.9cm}
%begin
{{\begin {scriptsize}\par
{ \center  Київський національний університет імені Тараса Шевченка \par}
{\parbox {5 cm}{Освітня програма \hfill}{\parbox {7 cm}{Системний аналіз\hfill}}\parbox {6 cm}{\hfill Семестр 2}}\par
{\parbox {5 cm} {Навчальний предмет \hfill}\parbox {7 cm}{Програмування \hfill}\parbox {6cm}{\hfill}\par}\vspace{-15pt} \end{scriptsize}}{\center Екзаменаційний білет \No \textbf { 22 }\par}}
{%beginitem
1. Статичні поля, статичні методи, методи класів.%enditem
\par
\vspace{2mm}

%beginitem
2. Кожен рядок текстового файлу містить по два дійсних числа, розділені білими символами.
Написати функцію, що для заданого текстового файлу будує новий текстовий файл, рядки якого крім зображень дйсних чисел ще містять результат обчислення на них функції $f$. У вихідному файлі зображення чисел мають розділятися пробілами. 
$$f(x, y) =\frac{\pi x + 14}{(\cos x - 0.7) (y - 43)}$$ 
Якість коду також оцінюється.

%enditem
\par
\vspace{2mm}

%beginitem
3. Реалізувати операцію перетину множин з повтореннями. Множина з повтореннями задається переліком своїх елементів, записаних за неспаданням у текстовому файлі через білі символи. Результат у такому ж самому форматі (і теж за неспаданням) має бути записаний у текстовий файл. Питома функція має отримувати три шляхи до файлів. Розмір вхідних множин такий, що вони не можуть бути завантажені повністю в пам’ять програми.

%enditem
}
{\smallskip\smallskip \begin {scriptsize} Затверджено на засіданні кафедри теоретичної кібернетики, протокол \No 4  від   25.01.2024 р.\par\smallskip\smallskip\smallskip\smallskip
{\parbox {6.5 cm} {Зав. кафедри \dotfill Ю.В.~Крак}}\hspace {0.7cm} {\hbox to 6.5 cm { Екзаменатор \dotfill Т.О.~Карнаух}} \end{scriptsize}}
%end
\newpage
%begin
{{\begin {scriptsize}\par
{ \center  Київський національний університет імені Тараса Шевченка \par}
{\parbox {5 cm}{Освітня програма \hfill}{\parbox {7 cm}{Системний аналіз\hfill}}\parbox {6 cm}{\hfill Семестр 2}}\par
{\parbox {5 cm} {Навчальний предмет \hfill}\parbox {7 cm}{Програмування \hfill}\parbox {6cm}{\hfill}\par}\vspace{-15pt} \end{scriptsize}}{\center Екзаменаційний білет \No \textbf { 23 }\par}}
{%beginitem
1. Алгоритм пошуку в глибину та задачі, які він може вирішувати.%enditem
\par
\vspace{2mm}

%beginitem
2. Кожен рядок текстового файлу містить по два дійсних числа, розділені білими символами.
Написати функцію, що для заданого текстового файлу будує новий текстовий файл, рядки якого крім зображень дйсних чисел ще містять результат обчислення на них функції $f$. У вихідному файлі зображення чисел мають розділятися пробілами. 
$$f(x, y) =\frac{\log_{10} x - 93}{(\tg x + 0.5) (y + 77)}$$ 
Якість коду також оцінюється.

%enditem
\par
\vspace{2mm}

%beginitem
3. Текстовий файл містить послідовність цілих чисел $a_1$, $a_2$, $\ldots$, $a_t$, записану у форматі одне число на рядок. Знайти сімнадцяте найменше значення, що записане у файлі. Кожне значення, що записане у файлі, враховується стільки разів, скільки входить у файл. Кількість чисел у файлі така, що вони не можуть бути завантажені повністю в пам'ять програми. Розв'язання оформити у вигляді функції.

%enditem
}
{\smallskip\smallskip \begin {scriptsize} Затверджено на засіданні кафедри теоретичної кібернетики, протокол \No 4  від   25.01.2024 р.\par\smallskip\smallskip\smallskip\smallskip
{\parbox {6.5 cm} {Зав. кафедри \dotfill Ю.В.~Крак}}\hspace {0.7cm} {\hbox to 6.5 cm { Екзаменатор \dotfill Т.О.~Карнаух}} \end{scriptsize}}
%end
\vspace{0.9cm}
%begin
{{\begin {scriptsize}\par
{ \center  Київський національний університет імені Тараса Шевченка \par}
{\parbox {5 cm}{Освітня програма \hfill}{\parbox {7 cm}{Системний аналіз\hfill}}\parbox {6 cm}{\hfill Семестр 2}}\par
{\parbox {5 cm} {Навчальний предмет \hfill}\parbox {7 cm}{Програмування \hfill}\parbox {6cm}{\hfill}\par}\vspace{-15pt} \end{scriptsize}}{\center Екзаменаційний білет \No \textbf { 24 }\par}}
{%beginitem
1. Багатовимірні масиви бібліотеки numpy: у чому відмінність моделювання матриці багатовимірним масивом бібліотеки numpy та як список списків.%enditem
\par
\vspace{2mm}

%beginitem
2. Кожен рядок текстового файлу містить по два дійсних числа, розділені білими символами.
Написати функцію, що для заданого текстового файлу будує новий текстовий файл, рядки якого крім зображень дйсних чисел ще містять результат обчислення на них функції $f$. У вихідному файлі зображення чисел мають розділятися пробілами. 
$$f(x, y) =\frac{\ln x + 91}{(\ctg x - 0.3) (y - 50)}$$ 
Якість коду також оцінюється.

%enditem
\par
\vspace{2mm}

%beginitem
3. Реалізувати операцію різниці множин з повтореннями. Множина з повтореннями задається переліком своїх елементів, записаних за неспаданням у текстовому файлі через білі символи. Результат у такому ж самому форматі (і теж за неспаданням) має бути записаний у текстовий файл. Питома функція має отримувати три шляхи до файлів. Розмір вхідних множин такий, що вони не можуть бути завантажені повністю в пам’ять програми.

%enditem
}
{\smallskip\smallskip \begin {scriptsize} Затверджено на засіданні кафедри теоретичної кібернетики, протокол \No 4  від   25.01.2024 р.\par\smallskip\smallskip\smallskip\smallskip
{\parbox {6.5 cm} {Зав. кафедри \dotfill Ю.В.~Крак}}\hspace {0.7cm} {\hbox to 6.5 cm { Екзаменатор \dotfill Т.О.~Карнаух}} \end{scriptsize}}
%end
\newpage
%begin
{{\begin {scriptsize}\par
{ \center  Київський національний університет імені Тараса Шевченка \par}
{\parbox {5 cm}{Освітня програма \hfill}{\parbox {7 cm}{Системний аналіз\hfill}}\parbox {6 cm}{\hfill Семестр 2}}\par
{\parbox {5 cm} {Навчальний предмет \hfill}\parbox {7 cm}{Програмування \hfill}\parbox {6cm}{\hfill}\par}\vspace{-15pt} \end{scriptsize}}{\center Екзаменаційний білет \No \textbf { 25 }\par}}
{%beginitem
1. Множини.%enditem
\par
\vspace{2mm}

%beginitem
2. Кожен рядок текстового файлу містить по два дійсних числа, розділені білими символами.
Написати функцію, що для заданого текстового файлу будує новий текстовий файл, рядки якого крім зображень дйсних чисел ще містять результат обчислення на них функції $f$. У вихідному файлі зображення чисел мають розділятися пробілами. 
$$f(x, y) =\frac{\log_{2} x + 39}{(\arccos x + 0.8) (y + 75)}$$ 
Якість коду також оцінюється.

%enditem
\par
\vspace{2mm}

%beginitem
3. Текстовий файл містить послідовність цілих чисел $a_1$, $a_2$, $\ldots$, $a_t$, записану у форматі одне число на рядок. Знайти тринадцяте найменше значення, що записане у файлі. Кожне значення, що записане у файлі, враховується тільки один раз. Кількість чисел у файлі така, що вони не можуть бути завантажені повністю в пам'ять програми. Розв'язання оформити у вигляді функції.

%enditem
}
{\smallskip\smallskip \begin {scriptsize} Затверджено на засіданні кафедри теоретичної кібернетики, протокол \No 4  від   25.01.2024 р.\par\smallskip\smallskip\smallskip\smallskip
{\parbox {6.5 cm} {Зав. кафедри \dotfill Ю.В.~Крак}}\hspace {0.7cm} {\hbox to 6.5 cm { Екзаменатор \dotfill Т.О.~Карнаух}} \end{scriptsize}}
%end
\vspace{0.9cm}
%begin
{{\begin {scriptsize}\par
{ \center  Київський національний університет імені Тараса Шевченка \par}
{\parbox {5 cm}{Освітня програма \hfill}{\parbox {7 cm}{Системний аналіз\hfill}}\parbox {6 cm}{\hfill Семестр 2}}\par
{\parbox {5 cm} {Навчальний предмет \hfill}\parbox {7 cm}{Програмування \hfill}\parbox {6cm}{\hfill}\par}\vspace{-15pt} \end{scriptsize}}{\center Екзаменаційний білет \No \textbf { 26 }\par}}
{%beginitem
1. Алгоритм пошуку в ширину та задачі, які він може вирішувати.%enditem
\par
\vspace{2mm}

%beginitem
2. Кожен рядок текстового файлу містить по два дійсних числа, розділені білими символами.
Написати функцію, що для заданого текстового файлу будує новий текстовий файл, рядки якого крім зображень дйсних чисел ще містять результат обчислення на них функції $f$. У вихідному файлі зображення чисел мають розділятися пробілами. 
$$f(x, y) =\frac{\sqrt x - 26}{(\arctg x - 0.2) (y - 76)}$$ 
Якість коду також оцінюється.

%enditem
\par
\vspace{2mm}

%beginitem
3. Кожен з двох текстових файлів містить послідовність дійсних чисел, записану у форматі одне число на рядок.  Перевірити, чи задають вони одну ту саму послідовність. Кількість чисел у кожному з двох файлів така, що вони не можуть бути завантажені повністю в пам'ять програми. Розв'язання оформити у вигляді функції.

%enditem
}
{\smallskip\smallskip \begin {scriptsize} Затверджено на засіданні кафедри теоретичної кібернетики, протокол \No 4  від   25.01.2024 р.\par\smallskip\smallskip\smallskip\smallskip
{\parbox {6.5 cm} {Зав. кафедри \dotfill Ю.В.~Крак}}\hspace {0.7cm} {\hbox to 6.5 cm { Екзаменатор \dotfill Т.О.~Карнаух}} \end{scriptsize}}
%end
\newpage
%begin
{{\begin {scriptsize}\par
{ \center  Київський національний університет імені Тараса Шевченка \par}
{\parbox {5 cm}{Освітня програма \hfill}{\parbox {7 cm}{Системний аналіз\hfill}}\parbox {6 cm}{\hfill Семестр 2}}\par
{\parbox {5 cm} {Навчальний предмет \hfill}\parbox {7 cm}{Програмування \hfill}\parbox {6cm}{\hfill}\par}\vspace{-15pt} \end{scriptsize}}{\center Екзаменаційний білет \No \textbf { 27 }\par}}
{%beginitem
1. Лінійні зв’язані структури даних.%enditem
\par
\vspace{2mm}

%beginitem
2. Кожен рядок текстового файлу містить по два дійсних числа, розділені білими символами.
Написати функцію, що для заданого текстового файлу будує новий текстовий файл, рядки якого крім зображень дйсних чисел ще містять результат обчислення на них функції $f$. У вихідному файлі зображення чисел мають розділятися пробілами. 
$$f(x, y) =\frac{\pi x + 97}{(\sin x - 0.6) (y - 42)}$$ 
Якість коду також оцінюється.

%enditem
\par
\vspace{2mm}

%beginitem
3. Текстовий файл містить послідовність цілих чисел $a_1$, $a_2$, $\ldots$, $a_t$, записану у форматі одне число на рядок. Знайти тринадцяте найбільше значення, що записане у файлі. Кожне значення, що записане у файлі, враховується тільки один раз. Кількість чисел у файлі така, що вони не можуть бути завантажені повністю в пам'ять програми. Розв'язання оформити у вигляді функції.

%enditem
}
{\smallskip\smallskip \begin {scriptsize} Затверджено на засіданні кафедри теоретичної кібернетики, протокол \No 4  від   25.01.2024 р.\par\smallskip\smallskip\smallskip\smallskip
{\parbox {6.5 cm} {Зав. кафедри \dotfill Ю.В.~Крак}}\hspace {0.7cm} {\hbox to 6.5 cm { Екзаменатор \dotfill Т.О.~Карнаух}} \end{scriptsize}}
%end
\vspace{0.9cm}
%begin
{{\begin {scriptsize}\par
{ \center  Київський національний університет імені Тараса Шевченка \par}
{\parbox {5 cm}{Освітня програма \hfill}{\parbox {7 cm}{Системний аналіз\hfill}}\parbox {6 cm}{\hfill Семестр 2}}\par
{\parbox {5 cm} {Навчальний предмет \hfill}\parbox {7 cm}{Програмування \hfill}\parbox {6cm}{\hfill}\par}\vspace{-15pt} \end{scriptsize}}{\center Екзаменаційний білет \No \textbf { 28 }\par}}
{%beginitem
1. Емпіричний підхід до визначення складності обчислень.%enditem
\par
\vspace{2mm}

%beginitem
2. Кожен рядок текстового файлу містить по два дійсних числа, розділені білими символами.
Написати функцію, що для заданого текстового файлу будує новий текстовий файл, рядки якого крім зображень дйсних чисел ще містять результат обчислення на них функції $f$. У вихідному файлі зображення чисел мають розділятися пробілами. 
$$f(x, y) =\frac{\ x - 31}{(\arcsin x + 0.2) (y + 92)}$$ 
Якість коду також оцінюється.

%enditem
\par
\vspace{2mm}

%beginitem
3. Реалізувати операцію симетричної різниці множин з повтореннями. Множина з повтореннями задається переліком своїх елементів, записаних за неспаданням у текстовому файлі через білі символи. Результат у такому ж самому форматі (і теж за неспаданням) має бути записаний у текстовий файл. Питома функція має отримувати три шляхи до файлів. Розмір вхідних множин такий, що вони не можуть бути завантажені повністю в пам’ять програми.

%enditem
}
{\smallskip\smallskip \begin {scriptsize} Затверджено на засіданні кафедри теоретичної кібернетики, протокол \No 4  від   25.01.2024 р.\par\smallskip\smallskip\smallskip\smallskip
{\parbox {6.5 cm} {Зав. кафедри \dotfill Ю.В.~Крак}}\hspace {0.7cm} {\hbox to 6.5 cm { Екзаменатор \dotfill Т.О.~Карнаух}} \end{scriptsize}}
%end
\newpage
%begin
{{\begin {scriptsize}\par
{ \center  Київський національний університет імені Тараса Шевченка \par}
{\parbox {5 cm}{Освітня програма \hfill}{\parbox {7 cm}{Системний аналіз\hfill}}\parbox {6 cm}{\hfill Семестр 2}}\par
{\parbox {5 cm} {Навчальний предмет \hfill}\parbox {7 cm}{Програмування \hfill}\parbox {6cm}{\hfill}\par}\vspace{-15pt} \end{scriptsize}}{\center Екзаменаційний білет \No \textbf { 29 }\par}}
{%beginitem
1. Файли формату csv.%enditem
\par
\vspace{2mm}

%beginitem
2. Кожен рядок текстового файлу містить по два дійсних числа, розділені білими символами.
Написати функцію, що для заданого текстового файлу будує новий текстовий файл, рядки якого крім зображень дйсних чисел ще містять результат обчислення на них функції $f$. У вихідному файлі зображення чисел мають розділятися пробілами. 
$$f(x, y) =\frac{\ln x - 38}{(\tg x + 0.3) (y + 78)}$$ 
Якість коду також оцінюється.

%enditem
\par
\vspace{2mm}

%beginitem
3. Реалізувати операцію перетину множин з повтореннями. Множина з повтореннями задається переліком своїх елементів, записаних за неспаданням у текстовому файлі через білі символи. Результат у такому ж самому форматі (і теж за неспаданням) має бути записаний у текстовий файл. Питома функція має отримувати три шляхи до файлів. Розмір вхідних множин такий, що вони не можуть бути завантажені повністю в пам’ять програми.

%enditem
}
{\smallskip\smallskip \begin {scriptsize} Затверджено на засіданні кафедри теоретичної кібернетики, протокол \No 4  від   25.01.2024 р.\par\smallskip\smallskip\smallskip\smallskip
{\parbox {6.5 cm} {Зав. кафедри \dotfill Ю.В.~Крак}}\hspace {0.7cm} {\hbox to 6.5 cm { Екзаменатор \dotfill Т.О.~Карнаух}} \end{scriptsize}}
%end
\vspace{0.9cm}
%begin
{{\begin {scriptsize}\par
{ \center  Київський національний університет імені Тараса Шевченка \par}
{\parbox {5 cm}{Освітня програма \hfill}{\parbox {7 cm}{Системний аналіз\hfill}}\parbox {6 cm}{\hfill Семестр 2}}\par
{\parbox {5 cm} {Навчальний предмет \hfill}\parbox {7 cm}{Програмування \hfill}\parbox {6cm}{\hfill}\par}\vspace{-15pt} \end{scriptsize}}{\center Екзаменаційний білет \No \textbf { 30 }\par}}
{%beginitem
1. Поняття часової та просторової складності обчислень. Складність у середньому. Асимптотичні оцінки складності.%enditem
\par
\vspace{2mm}

%beginitem
2. Кожен рядок текстового файлу містить по два дійсних числа, розділені білими символами.
Написати функцію, що для заданого текстового файлу будує новий текстовий файл, рядки якого крім зображень дйсних чисел ще містять результат обчислення на них функції $f$. У вихідному файлі зображення чисел мають розділятися пробілами. 
$$f(x, y) =\frac{\log_{10} x + 34}{(\arccos x - 0.5) (y - 74)}$$ 
Якість коду також оцінюється.

%enditem
\par
\vspace{2mm}

%beginitem
3. Текстовий файл містить послідовність цілих чисел $a_1$, $a_2$, $\ldots$, $a_t$, записану у форматі одне число на рядок. Знайти сімнадцяте найбільше значення, що записане у файлі. Кожне значення, що записане у файлі, враховується стільки разів, скільки входить у файл. Кількість чисел у файлі така, що вони не можуть бути завантажені повністю в пам'ять програми. Розв'язання оформити у вигляді функції.

%enditem
}
{\smallskip\smallskip \begin {scriptsize} Затверджено на засіданні кафедри теоретичної кібернетики, протокол \No 4  від   25.01.2024 р.\par\smallskip\smallskip\smallskip\smallskip
{\parbox {6.5 cm} {Зав. кафедри \dotfill Ю.В.~Крак}}\hspace {0.7cm} {\hbox to 6.5 cm { Екзаменатор \dotfill Т.О.~Карнаух}} \end{scriptsize}}
%end
\newpage
%begin
{{\begin {scriptsize}\par
{ \center  Київський національний університет імені Тараса Шевченка \par}
{\parbox {5 cm}{Освітня програма \hfill}{\parbox {7 cm}{Системний аналіз\hfill}}\parbox {6 cm}{\hfill Семестр 2}}\par
{\parbox {5 cm} {Навчальний предмет \hfill}\parbox {7 cm}{Програмування \hfill}\parbox {6cm}{\hfill}\par}\vspace{-15pt} \end{scriptsize}}{\center Екзаменаційний білет \No \textbf { 31 }\par}}
{%beginitem
1. Дерева. Обходи дерев.%enditem
\par
\vspace{2mm}

%beginitem
2. Кожен рядок текстового файлу містить по два дійсних числа, розділені білими символами.
Написати функцію, що для заданого текстового файлу будує новий текстовий файл, рядки якого крім зображень дйсних чисел ще містять результат обчислення на них функції $f$. У вихідному файлі зображення чисел мають розділятися пробілами. 
$$f(x, y) =\frac{\log_{10} x - 18}{(\ctg x + 0.9) (y - 29)}$$ 
Якість коду також оцінюється.

%enditem
\par
\vspace{2mm}

%beginitem
3. Текстовий файл містить послідовність цілих чисел $a_1$, $a_2$, $\ldots$, $a_t$, записану у форматі одне число на рядок. Знайти тринадцяте найменше значення, що записане у файлі. Кожне значення, що записане у файлі, враховується тільки один раз. Кількість чисел у файлі така, що вони не можуть бути завантажені повністю в пам'ять програми. Розв'язання оформити у вигляді функції.

%enditem
}
{\smallskip\smallskip \begin {scriptsize} Затверджено на засіданні кафедри теоретичної кібернетики, протокол \No 4  від   25.01.2024 р.\par\smallskip\smallskip\smallskip\smallskip
{\parbox {6.5 cm} {Зав. кафедри \dotfill Ю.В.~Крак}}\hspace {0.7cm} {\hbox to 6.5 cm { Екзаменатор \dotfill Т.О.~Карнаух}} \end{scriptsize}}
%end
\vspace{0.9cm}
%begin
{{\begin {scriptsize}\par
{ \center  Київський національний університет імені Тараса Шевченка \par}
{\parbox {5 cm}{Освітня програма \hfill}{\parbox {7 cm}{Системний аналіз\hfill}}\parbox {6 cm}{\hfill Семестр 2}}\par
{\parbox {5 cm} {Навчальний предмет \hfill}\parbox {7 cm}{Програмування \hfill}\parbox {6cm}{\hfill}\par}\vspace{-15pt} \end{scriptsize}}{\center Екзаменаційний білет \No \textbf { 32 }\par}}
{%beginitem
1. Словники.%enditem
\par
\vspace{2mm}

%beginitem
2. Кожен рядок текстового файлу містить по два дійсних числа, розділені білими символами.
Написати функцію, що для заданого текстового файлу будує новий текстовий файл, рядки якого крім зображень дйсних чисел ще містять результат обчислення на них функції $f$. У вихідному файлі зображення чисел мають розділятися пробілами. 
$$f(x, y) =\frac{\sqrt x + 21}{(\sin x - 0.7) (y + 67)}$$ 
Якість коду також оцінюється.

%enditem
\par
\vspace{2mm}

%beginitem
3. Текстовий файл містить послідовність цілих чисел $a_1$, $a_2$, $\ldots$, $a_t$, записану у форматі одне число на рядок. Знайти тринадцяте найбільше значення, що записане у файлі. Кожне значення, що записане у файлі, враховується тільки один раз. Кількість чисел у файлі така, що вони не можуть бути завантажені повністю в пам'ять програми. Розв'язання оформити у вигляді функції.

%enditem
}
{\smallskip\smallskip \begin {scriptsize} Затверджено на засіданні кафедри теоретичної кібернетики, протокол \No 4  від   25.01.2024 р.\par\smallskip\smallskip\smallskip\smallskip
{\parbox {6.5 cm} {Зав. кафедри \dotfill Ю.В.~Крак}}\hspace {0.7cm} {\hbox to 6.5 cm { Екзаменатор \dotfill Т.О.~Карнаух}} \end{scriptsize}}
%end
\newpage
\end{document}
